%% =====================================================================
%%  第6章 結論
%%  書くこと：本研究のまとめ（目的・方法・結果）と今後の課題
%%  本文に書いていない新しい結果をここで初めて出さない．
%% =====================================================================
\chapter{結論}
\label{chap:conclusion}

\section{まとめ}
\label{sec:summary}

\begin{writingnote}
「はじめに」の章で述べた目的と課題を振り返り，それに対して何を行い，何がわかったかを要約する．目的と課題は「はじめに」の章と同じ言葉で書くと，読者が対応を追いやすい．
主な結果は条件や数値を挙げて具体的に述べ，目的がどこまで達成できたかを明確にする．
そのうえで，「関連研究」の章で整理した既存研究に対して，本研究が何を新たに加えたのかを述べる．
数値実験で確かめた条件を超えて結論を一般化しない．
\end{writingnote}

本研究では，……を目的として，……を提案した．
数値実験の結果，提案手法は……であることを確認した．

\section{今後の課題}
\label{sec:future-work}

\begin{writingnote}
今後の課題は，本研究の限界のうち，今後の研究で取り組めるものから導く．限界は「提案手法」や「数値実験」の章で述べたものを取り上げ，この節で初めて持ち出さない．
「さらなる改良が必要である」のような漠然とした書き方ではなく，何が未解決で，どのように取り組むかを具体的に書く．
\end{writingnote}

%% まだ書いていない所には \TODO{...} を置いておくとよい．
%% 例：\TODO{今後の課題を具体的に書く}
%% review オプションを付けると赤字で表示され，付けないと表示されずに警告が出る．
今後の課題として，……が挙げられる．
